\section{Late-time Fourier identification and a pointwise sampling bound}
\label{sec:fourier}

Throughout this section $\lambda>0$ and $\sigma\geq0$ are constants,
$b=\lambda m$, and the selfsimilar daughter measure in
\eqref{eq:selfsimilar-daughters} is a fixed measure $B$. Thus
$B((0,1))=2$ and $\int\theta B(\dd\theta)=1$.
We assume a solution in \cref{def:solution}; no initial or daughter
logarithmic moment is required. The well-posed class in
\cref{thm:existence} supplies such solutions when $M_2(0)<\infty$.
Define the log-size characteristic function and fragmentation exponent by
\begin{equation}
 \chi_t(k)=\int_E e^{ik\log x}\eta_t(\dd x),\qquad
 \psi(k)=\sigma\int_{(0,1)}(\theta^{ik}-1)B(\dd\theta),
 \quad k\in\R.
 \label{eq:fourier-def}
\end{equation}
The symbol $\chi_t$ is a log-size characteristic function, not a
raw-size Laplace transform.
The exponent $\psi$ is continuous, vanishes at zero, and has nonpositive
real part. Its finite jump measure is
$\nu=\sigma(\log)_\# B$, supported on $(-\infty,0)$.

Fourier and Mellin inversion are established tools for fragmentation
inference. The exact pure-fragmentation transform evolution appears in
\citet[equations (12)--(14)]{DoumicEscobedo2016}; inverse recovery from
asymptotic profiles and from short-time observations is developed by
\citet{DoumicEscobedoTournus2018,DoumicEscobedoTournus2024}.
Most directly, \citet[Section 2.2]{Garnier2024} cancels an unknown factor
with two-time characteristic-function quotients and uses a distinguished
logarithm for decompounding. Here the task is to justify this cancellation
asymptotically in a nonlinear equation and bound the coagulation error.
The limiting factor need not be a characteristic function.
\citet[Sections 3.1.1--3.1.4]{HoangEtAl2022} already develops
log-size Fourier estimation from independent samples of an asymptotic
growth--fragmentation profile, including denominator regularization.
That model has size growth and division without coagulation.

\subsection{Factorization with an exponentially small nonlinear error}

Set, consistently with \eqref{eq:constant-log-parameters},
\begin{equation}
 \omega=\kappa(b+\sigma)>0,\qquad
 a_0=\frac{\lambda H(0)^2}{N_0},\qquad
 \delta(k)=\omega+\operatorname{Re}\psi(k),\qquad
 \Xi(k)=\frac{|k|a_0}{\delta(k)}\quad(\delta(k)>0).
 \label{eq:fourier-constants}
\end{equation}
The condition $\delta(k)>0$ defines an admissible frequency set containing
a neighborhood of zero; it can have several connected components.
Since $1-\cos u\leq u^2/2$,
\begin{equation}
 \operatorname{Re}\psi(k)
 \geq-\frac{\sigma k^2}2\int_{(0,1)}(\log\theta)^2B(\dd\theta),
 \label{eq:window-quadratic}
\end{equation}
so if $\sigma>0$ and the fraction law has a finite second logarithmic
moment, the set $\{\delta>0\}$ contains
\begin{equation}
 |k|<\left(\frac{2\omega}{\sigma\int(\log\theta)^2B(\dd\theta)}\right)^{1/2}.
 \label{eq:window}
\end{equation}
When $\sigma=0$, the set is all of $\R$. At criticality with equal
splitting, $\sigma=b$, $\omega=2\kappa b$, and
$\int(\log\theta)^2B(\dd\theta)=2(\log2)^2$, so \eqref{eq:window} reads
$|k|<\sqrt{2\kappa}/\log2\approx0.85$. In this case
$\operatorname{Re}\psi(k)=2b[\cos(k\log2)-1]$. With
$\theta_0=\arccos(1-\kappa)$, the full admissible set is therefore
\[
 \{k:\delta(k)>0\}
 =\bigcup_{j\in\mathbb Z}
 \left(\frac{2\pi j-\theta_0}{\log2},
       \frac{2\pi j+\theta_0}{\log2}\right).
\]
Its component containing zero is $|k|<\theta_0/\log2\approx0.86$.
These bands certify factorization; the nonvanishing amplitude needed for
quotient identification is guaranteed below only on a neighborhood of zero.
The admissible set is inherited from the rate $\omega$ in
\eqref{eq:half-decay}, whose asymptotic optimality is not asserted.
The factorization may hold at more frequencies, but it is proved only
where $\delta>0$.

\begin{theorem}[Nonlinear Fourier factorization]\label{thm:factorization}
For every real $k$ with $\delta(k)>0$, there is an amplitude
$C_\infty(k)$ such that, for all $t\geq0$,
\begin{align}
 |\chi_t(k)-e^{t\psi(k)}C_\infty(k)|
 &\leq \Xi(k)e^{-\omega t},\label{eq:factorization}\\
 |e^{-t\psi(k)}\chi_t(k)-C_\infty(k)|
 &\leq \Xi(k)e^{-\delta(k)t}.\label{eq:amplitude-error}
\end{align}
The set $\{\delta>0\}$ is open and contains zero, and $C_\infty$ is
continuous on it. There is a compact interval $I$ with zero in its
interior on which $\delta$ and $|C_\infty|$ are bounded below by positive
constants. Also $C_\infty(0)=1$.
\end{theorem}

\begin{proof}
Use the real and imaginary parts of $e^{ikz}$ in
\eqref{eq:normalized-log}. In the locally absolutely continuous sense,
\begin{align}
 \partial_t\chi_t(k)&=\psi(k)\chi_t(k)+r_t(k),\notag\\
 r_t(k)&=\frac\lambda{N(t)}\iint x
  [e^{ik\log(x+y)}-e^{ik\log x}]n_t(\dd x)n_t(\dd y).
 \label{eq:fourier-source}
\end{align}
The paired increment bound
$|e^{iku}-e^{ikv}|\leq|k||u-v|$ and
$\log(1+y/x)\leq\sqrt{y/x}$ give, by \eqref{eq:half-decay},
\begin{equation}
 |r_t(k)|\leq |k|\frac{\lambda H(t)^2}{N(t)}
             \leq |k|a_0e^{-\omega t}.
 \label{eq:fourier-forcing}
\end{equation}
Variation of constants therefore defines the absolutely convergent integral
\begin{equation}
 C_\infty(k)=\chi_0(k)+\int_0^\infty e^{-s\psi(k)}r_s(k)\dd s.
 \label{eq:fourier-amplitude}
\end{equation}
Subtracting its tail proves \eqref{eq:amplitude-error}; multiplying by
$e^{t\operatorname{Re}\psi(k)}$ proves \eqref{eq:factorization}.
Since $\psi$ is continuous with $\psi(0)=0$, the set $\{\delta>0\}$ is
open and contains zero. On each compact subset of it, $\delta$ is bounded
below by a positive constant and $\Xi$ is bounded, so
$e^{-t\psi}\chi_t$, which is continuous in $k$, converges uniformly to
$C_\infty$ by \eqref{eq:amplitude-error}. Thus the amplitude is
continuous on $\{\delta>0\}$. At zero it equals one, so a compact
neighborhood of zero gives the nonvanishing.
\end{proof}

The amplitude includes both the unknown initial law and the accumulated
nonlinear interactions. Equation \eqref{eq:fourier-amplitude} gives no
independence interpretation for it and does not assert positive
definiteness. The notation distinguishes it from the controlled clock
$C(t)$ in \eqref{eq:clocks}.

\subsection{Exact-data identification from local frequencies}

\begin{corollary}[Late-time quotients]\label{cor:fourier-ratio}
Let $I$ be the interval of \cref{thm:factorization}. Then
$\chi_t(k)\ne0$ for all sufficiently large $t$, uniformly in $k\in I$,
and for every $h>0$
\begin{equation}
 R_t(k):=\frac{\chi_{t+h}(k)}{\chi_t(k)}
       \longrightarrow e^{h\psi(k)}
 \label{eq:ratio-limit}
\end{equation}
uniformly on $I$. Whenever
$\Xi(k)e^{-\delta(k)t}\leq|C_\infty(k)|/2$,
\begin{equation}
 |R_t(k)-e^{h\psi(k)}|
 \leq \frac{2\Xi(k)}{|C_\infty(k)|}
 e^{h\operatorname{Re}\psi(k)}(1+e^{-\delta(k)h})e^{-\delta(k)t}.
 \label{eq:ratio-bias}
\end{equation}
Fix $h>0$. There is a compact interval $I'\subset I$ with zero in its
interior on which the continuous logarithm of the limiting quotient,
anchored at zero by $\log 1=0$, recovers $\psi$:
\begin{equation}
 \psi(k)=\frac1h\log\!\left(\lim_{t\to\infty}R_t(k)\right),
 \qquad k\in I'.
 \label{eq:local-exponent}
\end{equation}
\end{corollary}

\begin{proof}
Write $e^{-t\psi(k)}\chi_t(k)=C_\infty(k)+\epsilon_t(k)$.
The denominator has modulus at least $|C_\infty(k)|/2$ under the
stated condition. Subtract $e^{h\psi(k)}$ from the quotient and bound
$|\epsilon_{t+h}(k)-\epsilon_t(k)|$ using
\eqref{eq:amplitude-error}. Uniform bounds on $\Xi$, $\delta^{-1}$,
and $|C_\infty|^{-1}$ on $I$ give the uniform conclusion.
Choose $I'\subset I$ so small that $e^{h\psi(k)}$ stays near one and
$|h\operatorname{Im}\psi(k)|<\pi$ on $I'$. The principal logarithm there
is continuous and equals $h\psi(k)$; equivalently it is the unique
continuous branch anchored at zero.
\end{proof}

The interval $I$ depends on the kinetics and the initial data through
$\omega$, $a_0$, and the lower bound for $|C_\infty|$; the interval $I'$
depends in addition on $h$. These quantities are unknown in an
identification problem, and no computable procedure for choosing $I$ or
$I'$ from observations is offered here. The corollary is a proof of
principle.

For completeness, the analytic uniqueness ingredient needs no log moments.
It is the classical uniqueness theorem for boundary values of bounded
holomorphic functions in a half-plane, in a Paley--Wiener-type form: the
transform of a finite measure on a half-line extends to a bounded
holomorphic function in a half-plane, and a bounded holomorphic function
whose boundary values vanish on a set of positive length is identically
zero (a theorem of F.\ and M.\ Riesz, extended by Lusin and Privalov).
Continuity up to the boundary allows the short reflection proof below.

\begin{lemma}[One-sided transform uniqueness]\label{lem:one-sided}
Let $\nu_1,\nu_2$ be finite measures on $(-\infty,0)$. If
\[
 \int(e^{iky}-1)\nu_1(\dd y)
 =\int(e^{iky}-1)\nu_2(\dd y)
\]
on a nonempty open real interval, then $\nu_1=\nu_2$.
\end{lemma}

\begin{proof}
Put $\Delta=\nu_1-\nu_2$ and
$\mu=\Delta-\Delta(\R)\delta_0$. This is a finite signed measure
supported on $(-\infty,0]$, and its Fourier transform vanishes on the
specified interval. The function
\[
 f(z)=\int e^{izy}\mu(\dd y),\qquad \operatorname{Im}z<0,
\]
is holomorphic in the lower half-plane and continuous up to the real
boundary. Indeed, the exponential is bounded by one there; on compact
subsets of the open half-plane its damping controls every power of
$|y|$, justifying differentiation without moments.
Its zero boundary values are real. Schwarz reflection across a smaller
segment extends $f$ holomorphically through that segment. The extension
vanishes on the segment, so the identity theorem makes it zero on an
open subset and then throughout the lower half-plane. Boundary
continuity gives a zero Fourier transform on all of $\R$. Uniqueness
of Fourier transforms of finite signed measures yields $\mu=0$.
Restriction to $(-\infty,0)$ gives $\Delta=0$.
\end{proof}

\begin{theorem}[Structural identification]\label{thm:fourier-identification}
Ideal exact knowledge of the normalized size laws at the pairs of times
$(t_j,t_j+h)$, for one fixed lag $h>0$ and some sequence $t_j\to\infty$,
determines $\nu=\sigma(\log)_\#B$ uniquely.
If $\sigma>0$, it determines
\begin{equation}
 \sigma=\tfrac12\nu((-\infty,0)),\qquad
 B=\sigma^{-1}(\exp)_\#\nu.
 \label{eq:daughter-recovery}
\end{equation}
If $\sigma=0$, it identifies the zero selection rate but cannot identify
an unused $B$. With additionally known conserved mass $m$ and observed
unnormalized count growth $r_N$, it determines
\begin{equation}
 r_N=\frac{\log N(t+h)-\log N(t)}h=\sigma-\lambda m,
 \qquad \lambda=\frac{\sigma-r_N}{m}.
 \label{eq:coag-recovery}
\end{equation}
Neither the initial distribution nor its characteristic function needs
to be known.
\end{theorem}

\begin{proof}
Since the limit in \eqref{eq:ratio-limit} exists, the observations along
$t_j$ give \eqref{eq:local-exponent} on the open neighborhood
$\operatorname{int}I'$ of zero. For two candidate models, intersect
their neighborhoods;
\cref{lem:one-sided} identifies their jump measures. Their common total
mass is $2\sigma$, proving \eqref{eq:daughter-recovery} and the zero-rate
case. The count identity \eqref{eq:count} proves
\eqref{eq:coag-recovery}.
\end{proof}

A fixed change of size unit multiplies both $\chi_t$ and $\chi_{t+h}$ by
the same phase, which cancels in their quotient. Unknown initial data
cancel through $C_\infty$. Absolute mass and unnormalized count growth
are separate observations in \eqref{eq:coag-recovery}. The theorem uses
a continuum of exact frequencies in the window $I'$; finitely many
Fourier values do not supply the one-sided uniqueness conclusion, and
the passage from $\psi$ on $I'$ to $\nu$ is an analytic continuation from
a short interval, whose ill-posedness is quantified below. The measure $B$
records expected daughters and does not specify additional correlations
or genealogical information absent from the population balance.
The strict support assumption excludes zero log jumps. If one instead
allows $B$ on $(0,1]$ with the same count and mass constraints, an atom
at one contributes nothing directly to the exponent. It is nevertheless
identified by those constraints. Indeed, the visible measure
$\nu_-:=\sigma(\log)_\#(B|_{(0,1)})$ is determined by
\cref{lem:one-sided}, and
\[
 \sigma=\int_{(-\infty,0)}(1-e^y)\nu_-(\dd y),\qquad
 z:=\sigma B(\{1\})=2\sigma-\nu_-(\R).
\]
The integral has a bounded integrand. For $\sigma>0$ these formulas
recover $B=\sigma^{-1}[(\exp)_\#\nu_-+z\delta_1]$.
Thus strict reduction is not necessary for this algebraic identification
when both daughter constraints are known. The forward solution framework
in this paper retains its stated support assumption.

The holomorphic exponent satisfies the mass identity
$\psi(-i)=\sigma\int(\theta-1)B(\dd\theta)=-\sigma$.
This is an identity for the uniquely determined exponent, not an
instruction to continue noisy observations to a complex argument.
Identification from the window rests on \cref{lem:one-sided}, that is,
on analytic continuation from a short interval, and such continuation is
exponentially ill-posed. Quantitatively, let $g=\psi_1-\psi_2$ for two
admissible models with the same $\sigma$. Then $g$ is holomorphic in the
open lower half-plane, continuous up to the real axis, and bounded there
by $8\sigma$. If $|g|\leq\varepsilon$ on $I'$, the two-constants theorem
gives
\[
 |g(z)|\leq\varepsilon^{\varpi(z)}(8\sigma)^{1-\varpi(z)},
\]
where $\varpi(z)$ is the harmonic measure of $I'$ seen from $z$ in the
lower half-plane. This exponent tends to zero as $z$ approaches any real
frequency outside $I'$, so the data on $I'$ control no real frequency
outside $I'$ without further assumptions. If in addition
$\int\theta^{-a}B(\dd\theta)<\infty$ for some $a>0$, so that both
exponents extend holomorphically to the strip $|\operatorname{Im}z|<a$,
the same theorem in the strip slit along $I'$ gives an exponent that
decays exponentially with the distance from $I'$. The resulting bound
is then informative only up to distances of order
$\log\log(1/\varepsilon)$ from $I'$. The norm discontinuity that follows
is one symptom of this ill-posedness.
In fact, take $p\in(0,1)\setminus\{1/2\}$ and
$B_p=\delta_p+\delta_{1-p}$. As $q\to p$ with the four daughter atoms
distinct, the exponents for $B_q$ converge uniformly on each bounded
real-frequency interval, while
$\norm{B_q-B_p}_{\mathrm{var}}=4$.
Thus inversion to the daughter measure in variation norm is not continuous
under this data norm, even with fixed positive $\sigma$ and both daughter
constraints. This does not exclude stability in weaker metrics or under
additional regularity assumptions, but any such stability must be proved
separately; nothing here supplies it.

\subsection{Independent sampling: a pointwise sampling bound and observation time}

Fix one frequency $k\in I$ and a lag $h>0$. At the predetermined times
$t$ and $t+h$, observe $n$ independent draws from each of the deterministic
continuum number laws $\eta_t$ and $\eta_{t+h}$, with the two samples
also independent. These are hypothetical observations of continuum laws,
not the interacting particles of \cref{sec:finite}, and not auxiliary
paths from \cref{sec:auxiliary}.
Write $U=\chi_t(k)$, $V=\chi_{t+h}(k)$, and let
$\widehat U,\widehat V$ be the corresponding empirical averages of
$e^{ik\log X}$. Put $\widehat R=\widehat V/\widehat U$ when defined.
For $0<\varsigma<1$, set
\begin{equation}
 \varepsilon_n=2\sqrt{\frac{\log(8/\varsigma)}n}.
 \label{eq:sampling-epsilon}
\end{equation}

\begin{proposition}[Pointwise sampling bound for the ratio]\label{prop:ratio-sampling}
With probability at least $1-\varsigma$, both empirical characteristic
function errors are at most $\varepsilon_n$. On this event, whenever
$|\widehat U|>\varepsilon_n$, the exact finite-time ratio is defined and
\begin{equation}
 |\widehat R-R_t(k)|
 \leq\frac{\varepsilon_n(1+|\widehat R|)}
                  {|\widehat U|-\varepsilon_n}.
 \label{eq:observable-certificate}
\end{equation}
Let $\delta=\delta(k)>0$ be as in \eqref{eq:fourier-constants}, so that
$\omega-\delta=-\operatorname{Re}\psi(k)\geq0$. Choose
$\underline C>0$, $\Upsilon<\infty$, and $t_0\geq0$ so that
for $t\geq t_0$,
\begin{equation}
 |U|\geq\tfrac{\underline C}2e^{-(\omega-\delta)t},\qquad
 |R_t(k)-e^{h\psi(k)}|\leq\Upsilon e^{-\delta t}.
 \label{eq:sampling-signal}
\end{equation}
Such constants exist by \cref{thm:factorization,cor:fourier-ratio}.
If additionally $\varepsilon_n\leq(\underline C/4)e^{-(\omega-\delta)t}$,
then on the same event $\widehat U\ne0$ and
\begin{equation}
 |\widehat R-e^{h\psi(k)}|
 \leq\Upsilon e^{-\delta t}
  +\frac{4(2+\Upsilon)}{\underline C}\varepsilon_n e^{(\omega-\delta)t}.
 \label{eq:bias-noise}
\end{equation}
For all sufficiently large $n$, the deterministic schedule
\begin{equation}
 t_n=\frac{\log(1/\varepsilon_n)}\omega
 \label{eq:observation-time}
\end{equation}
satisfies these hypotheses and gives
\begin{equation}
 |\widehat R-e^{h\psi(k)}|
 \leq\left[\Upsilon+\frac{4(2+\Upsilon)}{\underline C}\right]
              \varepsilon_n^{\delta/\omega}
 \label{eq:balanced-ratio}
\end{equation}
with probability at least $1-\varsigma$.
\end{proposition}

\begin{proof}
Hoeffding's bounded-variable inequality \citep{Hoeffding1963}, applied
to each real and imaginary component in $[-1,1]$ at threshold
$\varepsilon/\sqrt2$, and a union bound over all four components give
\[
 \Prob\bigl(|\widehat U-U|>\varepsilon
             \text{ or }|\widehat V-V|>\varepsilon\bigr)
 \leq8e^{-n\varepsilon^2/4}.
\]
Substitute \eqref{eq:sampling-epsilon}. On the resulting event,
$|U|\geq|\widehat U|-\varepsilon_n$, and
\[
 \widehat R-R_t
 =\frac{\widehat R(U-\widehat U)-(V-\widehat V)}U
\]
proves \eqref{eq:observable-certificate}. Under the stronger signal
condition, $|\widehat U|\geq(\underline C/4)e^{-(\omega-\delta)t}$ and
$|R_t|\leq1+\Upsilon$, because $|e^{h\psi(k)}|\leq1$.
Using instead
\[
 \widehat R-R_t
 =\frac{\widehat V-V-R_t(\widehat U-U)}{\widehat U}
\]
gives the second term in \eqref{eq:bias-noise}; add the deterministic
bias. Finally, at \eqref{eq:observation-time} both
$e^{-\delta t_n}$ and $\varepsilon_ne^{(\omega-\delta)t_n}$ equal
$\varepsilon_n^{\delta/\omega}$, so both terms are constant multiples of
this quantity. It tends to zero, which verifies the signal condition
eventually; also $t_n\to\infty$.
\end{proof}

The pointwise sampling bound \eqref{eq:observable-certificate} controls
sampling error to the exact finite-time ratio; the bias in
\eqref{eq:ratio-bias} must still be added to estimate $e^{h\psi(k)}$.
Its strict denominator condition is required. At fixed confidence,
\eqref{eq:balanced-ratio} is an upper bound of order
$n^{-\delta/(2\omega)}$, attained at a time of order
$(\log n)/(2\omega)$. This is a sufficient balance of the proved bounds,
not a minimax result or an optimal experimental-design theorem.
Both the schedule and the exponent use the certified decay coefficient
$\omega=\kappa(b+\sigma)$ for $H^2/N$ in \eqref{eq:half-decay}.
Its asymptotic optimality is not asserted. For the Golovin
benchmark of pure additive coagulation from monodisperse data
\citep{Golovin1963}, $H^2/N$ decays roughly like $e^{-bt}$ up to
polynomial factors, against the certified $\omega=\kappa b\approx0.17b$.
This benchmark shows that a stronger decay estimate can permit an earlier
observation time, but does not prove faster decay throughout the model
class. If a stronger exponential forcing bound is established for a given
model, the balance can be recomputed with that bound. At frequencies with
$\operatorname{Re}\psi(k)<0$, increasing the decay coefficient improves
the resulting sampling exponent $\delta/(2\omega)$.
When $\delta=\omega$, that is, when $\operatorname{Re}\psi(k)=0$, the
sampling term does not grow with time and the same schedule remains
sufficient. The sampling exponent is then already $1/2$, including in the
pure-coagulation benchmark, and a stronger decay bound does not increase
it. The constants depend on unknown kinetics
and initial data; using the schedule requires their relevant values or
prior class bounds. No adaptive choice of time or frequency is proved,
and no computable procedure selects the frequency interval $I$; the
proposition is a proof of principle under a hypothetical observation
model.
A local logarithm separated from its branch cut transfers pointwise ratio
convergence to the exponent, but supplies no full noisy daughter inversion.
The logarithmic observation time in sample size also does not validate
an estimator based on a single finite reactor: the continuum obstruction
in \cref{thm:finite-discrepancy} concerns a different population parameter
and a different observation law.
